\documentclass[12pt,a4paper]{article}

\usepackage{fontspec}
\setmainfont{Times New Roman}
\setmonofont{Courier New}
\newfontfamily\cyrillicfonttt{Courier New}
\usepackage{polyglossia}
\setdefaultlanguage{russian}
\usepackage[margin=20mm]{geometry}
\usepackage{amsmath,amssymb}
\usepackage{tikz}
\usetikzlibrary{positioning,arrows.meta}
\usepackage{xcolor}
\usepackage{listings}
\usepackage{longtable}
\usepackage{hyperref}

\definecolor{ink}{HTML}{20354A}
\definecolor{accent}{HTML}{286A8A}
\definecolor{softblue}{HTML}{E7F1F6}
\definecolor{softyellow}{HTML}{FFF4D6}
\definecolor{codebg}{HTML}{F4F6F8}

\hypersetup{colorlinks=true,linkcolor=accent,urlcolor=accent}
\setlength{\parindent}{1.2em}
\setlength{\parskip}{0.35em}
\renewcommand{\arraystretch}{1.15}

\lstdefinestyle{gostyle}{
  language=Go,
  basicstyle=\ttfamily\small,
  keywordstyle=\color{accent}\bfseries,
  commentstyle=\color{gray}\itshape,
  stringstyle=\color{red!60!black},
  backgroundcolor=\color{codebg},
  frame=single,
  rulecolor=\color{black!15},
  numbers=left,
  numberstyle=\tiny\color{gray},
  numbersep=7pt,
  breaklines=true,
  columns=fullflexible,
  keepspaces=true,
  showstringspaces=false,
  tabsize=4
}

\newcommand{\code}[1]{\texttt{#1}}
\newcommand{\sectiontitle}[1]{\section{\color{ink}#1}}

\begin{document}

\begin{center}
  {\LARGE\bfseries\color{ink}Клеточный автомат простыми словами}\\[0.4em]
  {\large Окрестность фон Неймана и разбор пакета \code{automation}}\\[0.6em]
  {Разбор проекта лабораторной работы, вариант 2}
\end{center}

\vspace{0.5em}
\noindent\colorbox{softblue}{\parbox{\dimexpr\linewidth-2\fboxsep}{
  \textbf{Идея пособия.} Понять, что происходит с клетками на одном шаге,
  как пять клеток превращаются в номер от 0 до 31 и как этот номер
  используется в Go-коде.
}}

\sectiontitle{Что такое клеточный автомат}

Клеточный автомат - это сетка клеток, которая меняется шаг за шагом.
Каждая клетка имеет одно из заранее заданных состояний. В нашей программе
состояний два: клетка либо живая, либо мёртвая.

Состояние клетки на следующем шаге зависит не от всего поля сразу, а только
от небольшой группы клеток рядом с ней. Это правило применяется одинаково
ко всем клеткам поля.

Можно представить поле как фотографию. Чтобы построить следующую фотографию,
программа смотрит на старое поле, вычисляет новое состояние каждой клетки,
записывает результаты отдельно и только потом заменяет старое поле новым.
Поэтому все клетки обновляются одновременно: результат для одной клетки
не влияет на расчёт другой клетки в том же шаге.

В коде состояние хранится как значение типа \code{bool}:
\code{true} - живая клетка, \code{false} - мёртвая.

\sectiontitle{Окрестность фон Неймана}

Окрестность - это клетки, чьи состояния учитываются при обновлении одной
выбранной клетки. В окрестности фон Неймана берутся четыре клетки,
соприкасающиеся с центральной общей стороной: сверху, снизу, слева и справа.
В нашем правиле учитывается также сама центральная клетка, поэтому всего
получается пять значений.

\begin{center}
\begin{tikzpicture}[
  box/.style={draw=ink, very thick, minimum width=2.2cm, minimum height=0.9cm,
              align=center, rounded corners=2pt, font=\small},
  every node/.style={font=\small}
]
  \node[box,fill=softblue] (center) at (0,0) {$s_0$\\центр};
  \node[box,fill=accent!14] (up) at (0,1.35) {$s_1$\\верх};
  \node[box,fill=accent!14] (right) at (2.8,0) {$s_4$\\право};
  \node[box,fill=accent!14] (down) at (0,-1.35) {$s_2$\\низ};
  \node[box,fill=accent!14] (left) at (-2.8,0) {$s_3$\\лево};
  \draw[-{stealth},thick,accent] (up) -- (center);
  \draw[-{stealth},thick,accent] (right) -- (center);
  \draw[-{stealth},thick,accent] (down) -- (center);
  \draw[-{stealth},thick,accent] (left) -- (center);
\end{tikzpicture}

Рисунок 1 - клетки, состояния которых используются для обновления центра
\end{center}

Диагональные клетки в эту окрестность не входят. Это отличает её, например,
от окрестности Мура, где учитываются также четыре диагонали.

В программе значения идут в строго заданном порядке:

\[
\text{центр},\quad \text{верх},\quad \text{низ},\quad
\text{лево},\quad \text{право}.
\]

Важно не менять этот порядок только в одном месте: по нему и формируется
номер конфигурации, и выбирается результат из правила.

\sectiontitle{Как пять состояний превращаются в номер}

У каждой из пяти клеток два возможных состояния: 0 или 1. Поэтому всего
есть $2^5=32$ разных комбинации. Каждой комбинации соответствует номер
от 0 до 31.

Пусть $a,b,c,d,e$ - состояния центра, верхней, нижней, левой и правой
клеток соответственно. Номер комбинации - это просто двоичное число,
составленное из этих пяти значений:

\[
k = 16a + 8b + 4c + 2d + e.
\]

Например, для набора $1,0,1,0,0$ получаем двоичную запись $10100_2$,
а значит номер $k=20$.

В \code{field.go} этот расчёт выполняет \code{neighbourhoodIndex}:

\begin{lstlisting}[style=gostyle]
bits := [5]bool{
    f.At(row, col),   // центр
    f.At(row-1, col), // верх
    f.At(row+1, col), // низ
    f.At(row, col-1), // лево
    f.At(row, col+1), // право
}
index := 0
for _, bit := range bits {
    index <<= 1
    if bit {
        index++
    }
}
\end{lstlisting}

Переменная \code{bits} содержит пять состояний. Цикл каждый раз сдвигает
накопленный номер влево на один разряд и, если очередное состояние равно
\code{true}, добавляет единицу.

Для набора $1,0,1,0,0$ это выглядит так:

\begin{center}
\begin{tabular}{c|ccccc}
  Прочитано & 1 & 0 & 1 & 0 & 0 \\
  Номер после шага & 1 & 2 & 5 & 10 & 20 \\
\end{tabular}
\end{center}

То есть цикл не ищет комбинацию в списке. Он сам собирает её номер из битов.

\sectiontitle{Что хранится в \code{Rule}}

В Go нет классов в привычном смысле, здесь используются структуры и методы.
Структура \code{Rule} хранит номер правила:

\begin{lstlisting}[style=gostyle]
type Rule struct {
    Number uint32
}
\end{lstlisting}

Тип \code{uint32} означает беззнаковое целое из 32 бит. Поскольку у пяти
клеток 32 комбинации, каждый бит числа хранит ответ для одной комбинации:
0 - центр должен стать мёртвым, 1 - живым.

Для варианта 2 в проекте используется формула:

\[
R = 2 \cdot 11 \cdot 2006 \cdot 13 \cdot 12
  = 6\,884\,592
  = 00000000011010010000110011110000_2.
\]

Нули слева добавлены, чтобы запись содержала ровно 32 бита. В методе
\code{Bits} формат \code{\%032b} как раз означает: вывести число в двоичной
системе и дополнить его нулями слева до 32 символов.

Метод \code{transition} проверяет нужный бит:

\begin{lstlisting}[style=gostyle]
func (r Rule) transition(configuration int) bool {
    return r.Number&(1<<uint(31-configuration)) != 0
}
\end{lstlisting}

Здесь \code{configuration} - номер комбинации от 0 до 31. Самый первый
набор, \code{00000}, соответствует старшему (самому левому) биту. Поэтому
для перехода к номеру $k$ код сдвигает единицу на $31-k$ позиций, оставляет
только этот бит операцией \code{\&} и проверяет, не равен ли результат нулю.

\noindent\colorbox{softyellow}{\parbox{\dimexpr\linewidth-2\fboxsep}{
  \textbf{Пример.} Набор $10100$ имеет номер 20. Для него проверяется бит
  с номером $31-20=11$ при счёте справа с нуля. В правиле варианта 2 этот
  бит равен 1, поэтому новая центральная клетка будет живой.
}}

\sectiontitle{Что хранится в \code{Field}}

\code{Field} хранит всё текущее поле и правило, по которому оно меняется:

\begin{lstlisting}[style=gostyle]
type Field struct {
    cells [][]bool
    rule  Rule
}
\end{lstlisting}

\begin{itemize}
  \item \code{cells} - квадратная матрица состояний. Первый индекс -
        строка, второй - столбец.
  \item \code{rule} - копия правила, которое применяется при каждом шаге.
\end{itemize}

Основные методы поля:

\begin{itemize}
  \item \code{New(size, rule)} создаёт квадратное поле. Новые значения
        \code{bool} изначально равны \code{false}, значит все клетки мёртвые.
  \item \code{Size()} возвращает размер стороны поля.
  \item \code{At(row, col)} читает состояние клетки.
  \item \code{Set(row, col, state)} задаёт клетке состояние.
  \item \code{Toggle(row, col)} меняет состояние на противоположное.
  \item \code{Clear()} делает все клетки мёртвыми.
  \item \code{Randomize(seed)} случайно заполняет поле. Одинаковый \code{seed}
        даёт одинаковую последовательность случайных состояний.
\end{itemize}

Чтение и изменение клетки проходят через \code{wrap}. Эта функция реализует
тороидальные границы: если выйти за край, координата появляется с другой
стороны поля. При размере 5 координата $-1$ превращается в 4, а координата
5 - в 0. В коде используется формула

\[
\operatorname{wrap}(x)=((x\%N)+N)\%N,
\]

где $N$ - размер поля. Добавление $N$ нужно, чтобы результат был
неотрицательным и оставался допустимым индексом от 0 до $N-1$.

\sectiontitle{Как работает один шаг}

Главный метод обновления - \code{Step}. Он проходит по каждой клетке,
находит номер её окрестности, получает ответ из \code{Rule} и записывает
его во временное поле:

\begin{lstlisting}[style=gostyle]
func (f *Field) Step() {
    next := makeCells(f.Size())
    for row := range f.cells {
        for col := range f.cells[row] {
            next[row][col] = f.rule.transition(
                f.neighbourhoodIndex(row, col),
            )
        }
    }
    f.cells = next
}
\end{lstlisting}

По порядку здесь происходит следующее:

\begin{enumerate}
  \item \code{makeCells} создаёт пустую матрицу для нового поколения.
  \item Вложенные циклы перебирают все строки и столбцы.
  \item \code{neighbourhoodIndex} читает центр и четыре соседние клетки,
        затем превращает их состояния в число 0--31.
  \item \code{transition} проверяет соответствующий бит правила.
  \item Результат записывается в \code{next}, а не в старое поле.
  \item Когда все клетки рассчитаны, строка \code{f.cells = next} заменяет
        старое поколение новым.
\end{enumerate}

Именно временная матрица обеспечивает одновременное обновление: каждая
клетка нового поколения рассчитывается по одному и тому же старому полю.
Для поля $N\times N$ один вызов \code{Step} обрабатывает $N^2$ клеток,
поэтому время работы растёт пропорционально площади поля.

\sectiontitle{Вся цепочка в одном месте}

\begin{center}
\begin{tikzpicture}[
  node distance=0.65cm,
  block/.style={draw=ink,rounded corners,fill=softblue,align=center,
                minimum width=3.3cm,minimum height=0.85cm,font=\small},
  arrow/.style={-{stealth},very thick,accent}
]
  \node[block] (step) {\code{Field.Step()}};
  \node[block,below=of step] (neighbors) {состояния центра\\и 4 соседей};
  \node[block,below=of neighbors] (index) {номер комбинации\\от 0 до 31};
  \node[block,below=of index] (rule) {\code{Rule.transition()}\\бит правила};
  \node[block,below=of rule] (next) {новое состояние\\в матрице \code{next}};
  \draw[arrow] (step) -- (neighbors);
  \draw[arrow] (neighbors) -- (index);
  \draw[arrow] (index) -- (rule);
  \draw[arrow] (rule) -- (next);
\end{tikzpicture}
\end{center}

Коротко: \textbf{смотрим на пять клеток → кодируем их нулями и единицами →
получаем номер → читаем бит правила → записываем новое состояние}.

\vfill
\begin{center}
  {\small Материал составлен по текущему коду проекта в
  \code{lab1/internal/automation}.}
\end{center}

\end{document}
